\chapter{Gaz naturel et GNL}\label{ch:m3:natural-gas-and-lng}

Au début de 2020, un mégawattheure de gaz livré le mois suivant aux Pays-Bas se négociait
nettement en dessous de 25~euros. Deux ans et demi plus tard, il dépassait 300~euros. Durant ces mêmes
mois, un méthanier pouvait charger du gaz en Louisiane à un prix lié à l'indice de référence américain,
soit quelques dollars par million de British thermal units, faire route pendant deux semaines
et vendre la cargaison en Europe à un prix plusieurs fois supérieur. Le gaz est une matière
première de réseau~: il circule dans des gazoducs et, là où ceux-ci s'arrêtent, à bord de
navires qui doivent d'abord le refroidir à $-\qty{162}{\degreeCelsius}$. Ses prix sont locaux, saisonniers
et, lorsque le réseau se rompt, extrêmes. Ce chapitre présente les hubs où le gaz est coté, la base
entre ces hubs, le cycle saisonnier du stockage, l'arbitrage des cargaisons liquéfiées et la crise
de 2022.

\section{Le gaz comme matière première de réseau~: hubs et unités}

Le gaz s'échange sur des hubs, c'est-à-dire des points d'un réseau de gazoducs auxquels la
propriété peut passer d'une partie à une autre. Certains hubs sont des lieux physiques~; en Europe, la plupart sont notionnels.

\begin{definition}[Point d'échange virtuel, Henry Hub, Title Transfer Facility]\label{def:m3:natural-gas-and-lng:hubs}
Un \emph{point d'échange virtuel}\index{point d'échange virtuel} est un point notionnel d'un
réseau de transport de gaz, exploité par son gestionnaire de réseau, auquel le gaz entré dans
le système peut être échangé sans égard à l'endroit où il se trouve physiquement.
Le \emph{Henry Hub}\index{Henry Hub} est une interconnexion physique de gazoducs en
Louisiane, dont le prix sert d'indice de référence américain et qui constitue le point de livraison des
principaux contrats à terme américains sur le gaz. Le \emph{Title Transfer Facility}\index{Title Transfer Facility} (TTF) est le
point d'échange virtuel du réseau néerlandais et le principal indice de référence européen du gaz.
\end{definition}

Les unités diffèrent selon les régions. Les États-Unis cotent le gaz en dollars par million de
British thermal units (MMBtu)~; l'Europe continentale en euros par mégawattheure d'énergie~; le
Royaume-Uni en pence par therm (\qty{0.1}{\mmbtu}). Un MMBtu vaut \qty{0.293071}{MWh}~:
un prix TTF de \qty{40}{EUR/MWh} avec un EURUSD à 1.16 correspond à \qty{13.60}{\$/\mmbtu}. Le
contrat à terme américain porte sur 10\,000 MMBtu livrés au \omterm{def:m3:natural-gas-and-lng:hubs}{Henry Hub}~; sa négociation prend fin
trois jours ouvrés avant le premier jour du mois de livraison.

\begin{definition}[Jour gazier]\label{def:m3:natural-gas-and-lng:gasday}
Le \emph{jour gazier}\index{jour gazier} est la période de 24~heures sur laquelle un réseau équilibre
les flux et procède au règlement-livraison des produits journaliers. Dans l'Union européenne, il court
de 05:00 à 05:00 UTC en hiver et de 04:00 à 04:00 UTC en été.
\end{definition}

\section{Base géographique et gazoducs}

Un prix en un hub est un prix du gaz en ce lieu. Entre deux hubs, l'écart de prix
est borné par le coût de transport du gaz entre eux lorsque la capacité est disponible, et
ne l'est plus lorsqu'elle est saturée.

\begin{definition}[Base géographique]\label{def:m3:natural-gas-and-lng:basis}
La \emph{base géographique}\index{base géographique} d'un hub est son prix moins le
prix d'un hub de référence (aux États-Unis, le \omterm{def:m3:natural-gas-and-lng:hubs}{Henry Hub}) pour la même période de livraison. Elle
s'échange sous forme de swaps de base et de contrats à terme.
\end{definition}

\begin{proposition}[La base n'est bornée par le transport que lorsque la capacité est disponible]\label{prop:m3:natural-gas-and-lng:basis}
Soit $c$ le coût variable de transport du gaz du hub $A$ vers le hub $B$, et supposons qu'un
expéditeur détienne une capacité inutilisée. Alors $P_B - P_A \le c$~: sinon l'expéditeur achète en
$A$, transporte et vend en $B$ pour un profit sans risque. Sans capacité inutilisée, l'écart
peut être de n'importe quelle taille, et sa valeur revient aux détenteurs de capacité.
\end{proposition}

\begin{proof}
Avec capacité, la transaction $-P_A - c + P_B > 0$ est accessible à quiconque la détient et
est réalisée jusqu'à ce que les prix bougent. Sans capacité, aucune transaction ne les relie.
\end{proof}

Une région productrice dont les gazoducs sortants sont saturés peut voir le prix de son hub chuter
loin en dessous de la référence, parfois sous zéro~; une région consommatrice coupée de
l'approvisionnement peut le voir grimper loin au-dessus. La valeur de la capacité de transport est
l'espérance de ces écarts, ce qui explique que la capacité elle-même soit vendue aux enchères.

\section{Stockage et saisonnalité}

La demande de gaz culmine en hiver (chauffage) et, dans les régions chaudes, en été (électricité
pour la climatisation), tandis que la production est à peu près constante. Le stockage fait le lien entre les deux.

\begin{definition}[Saison d'injection, spread été-hiver]\label{def:m3:natural-gas-and-lng:season}
La \emph{saison d'injection}\index{saison d'injection} est la période de l'année pendant laquelle
le stockage est rempli, environ d'avril à octobre dans l'hémisphère nord~; le gaz est
soutiré durant les mois d'hiver qui suivent. Le \emph{spread été-hiver}\index{spread été-hiver}
est le prix du gaz pour l'hiver suivant moins
le prix pour l'été précédent~; il rémunère le stockage du gaz de l'été à l'hiver.
\end{definition}

\begin{omfigure}
\centering
\begin{tikzpicture}[font=\small]
\begin{axis}[width=11cm,height=4.6cm,
  xlabel={mois de livraison},ylabel={prix à terme (stylisé)},
  tick label style={font=\small},label style={font=\small},
  xmin=0.5,xmax=12.5,ymin=0,ymax=20,grid=major,grid style={gray!25},
  xtick={1,2,3,4,5,6,7,8,9,10,11,12},
  xticklabels={avr.,mai,juin,juil.,août,sept.,oct.,nov.,déc.,janv.,févr.,mars},
  yticklabels={}]
\addplot[omDef,very thick,mark=*,mark size=1.4pt] coordinates {(1,8.6) (2,8.3) (3,8.2) (4,8.3) (5,8.4) (6,8.6) (7,9.0) (8,10.2) (9,11.2) (10,11.6) (11,11.3) (12,10.4)};
\draw[omThm,thick,{Stealth[length=2mm]}-{Stealth[length=2mm]}] (axis cs:3,8.2) -- (axis cs:3,11.6);
\node[omThm,font=\footnotesize,anchor=north west,align=left,fill=white,inner sep=1pt] at (axis cs:0.7,19.2) {spread été-hiver~: rémunère\\le stockage du gaz de l'été à l'hiver};
\draw[gray,dashed] (axis cs:3,11.6) -- (axis cs:10,11.6);
\node[font=\footnotesize,anchor=south] at (axis cs:4,1.2) {injection};
\node[font=\footnotesize,anchor=south] at (axis cs:10,1.2) {soutirage};
\end{axis}
\end{tikzpicture}

\omcaption{La forme saisonnière d'une courbe à terme du gaz, schématique~: les mois d'été sont
bon marché, les mois d'hiver chers. Un propriétaire de stockage achète l'été et vend l'hiver~; le spread,
net des coûts d'injection et de soutirage, constitue sa valeur intrinsèque. Stylisé, sans
données.}\label{fig:m3:natural-gas-and-lng:season}
\end{omfigure}

La valorisation la plus simple du stockage achète les mois d'été les moins chers et vend à terme les
mois d'hiver les plus chers, dans la limite de la capacité de l'installation et de ses taux d'injection
et de soutirage~; le reste de sa valeur tient à la réoptimisation au fil des mouvements de prix, que
livre~6, chapitre~16, traite comme des options.

\begin{dated}{2026-09}{Règles européennes de stockage}\label{dat:m3:natural-gas-and-lng:storage}
La Regulation (EU) 2022/1032 exigeait que les stocks des États membres soient remplis à 80~\% au
1\up{er} novembre 2022, puis à 90~\% au 1\up{er} novembre de chaque année suivante. La Regulation (EU) 2025/1733, publiée le
10 septembre 2025, prolonge les règles jusqu'à la fin de 2027 et permet que l'objectif de 90~\% soit atteint
à tout moment entre le 1\up{er} octobre et le 1\up{er} décembre, avec des écarts pouvant aller jusqu'à 10 points de pourcentage
dans des conditions de marché difficiles.
\end{dated}

\section{GNL~: liquéfaction, transport maritime et arbitrage de cargaisons}

\begin{definition}[Gaz naturel liquéfié]\label{def:m3:natural-gas-and-lng:lng}
\emph{gaz naturel liquéfié}\index{gaz naturel liquéfié} (GNL, LNG) est du gaz naturel refroidi
à environ $-\qty{162}{\degreeCelsius}$, température à laquelle il devient un liquide environ six cents
fois plus dense que le gaz~: il est transporté dans des méthaniers isolés, en perdant une petite part
par évaporation (boil-off) en cours de route, puis est réchauffé et reconverti en gaz (regazéifié) au terminal de réception.
\end{definition}

Le GNL transforme les marchés gaziers régionaux en un marché mondial, mais seulement à la marge, parce que la
chaîne est coûteuse (\cref{fig:m3:natural-gas-and-lng:chain}). Ses contrats se répartissent en deux
familles. Les contrats à long terme conclus avec les acheteurs asiatiques étaient longtemps indexés sur le pétrole.

\begin{definition}[Indexation sur le pétrole, flexibilité de destination]\label{def:m3:natural-gas-and-lng:contracts}
L'\emph{indexation sur le pétrole}\index{indexation sur le pétrole} fixe le prix du gaz par une formule $P = s \cdot P_{\mathrm{oil}}
+ k$, où la pente $s$ multiplie une moyenne des prix passés du pétrole, à laquelle s'ajoute une constante. Un contrat a une
\emph{flexibilité de destination}\index{flexibilité de destination} lorsque l'acheteur peut enlever la
cargaison vers n'importe quelle destination plutôt que vers un terminal désigné.
\end{definition}

Le modèle d'exportation américain est différent~: l'acheteur paie un droit fixe par MMBtu pour
disposer de capacités de liquéfaction, qu'il enlève ou non des cargaisons, plus un prix variable égal à 115~\% du
prix \omterm{def:m3:natural-gas-and-lng:hubs}{Henry Hub} par cargaison enlevée, \omterm{def:m3:physical-commodity-markets:incoterms}{franco à bord} (FOB) au terminal~; il peut annuler une cargaison
moyennant préavis, tout en restant redevable du droit fixe. Un tel acheteur détient un portefeuille d'options~: chaque mois, il enlève la cargaison
si le meilleur \omterm{def:m3:physical-commodity-markets:netback}{netback} couvre le prix variable, et l'envoie vers la destination qui
la rémunère le mieux.

\begin{definition}[Japan Korea Marker]\label{def:m3:natural-gas-and-lng:jkm}
Le \emph{Japan Korea Marker}\index{Japan Korea Marker} (JKM) est l'évaluation quotidienne établie par une \omterm{def:m3:physical-commodity-markets:pra}{agence de cotation des prix}
(PRA) des cargaisons spot de GNL livrées ex-ship au Japon, en Corée du Sud,
en Chine et à Taïwan, principal indice de référence spot asiatique.
\end{definition}

\begin{omfigure}
\centering
\begin{tikzpicture}[font=\small,
  box/.style={draw,thick,rounded corners=2pt,align=center,minimum height=1cm,text width=2.2cm},
  arr/.style={-{Stealth[length=2.2mm]},thick}]
\node[box,draw=omDef,fill=omDef!8] (hh) at (0,0) {gaz par gazoduc\\à Henry Hub};
\node[box,draw=omProp,fill=omProp!8] (liq) at (3.2,0) {liquéfaction\\(FOB)};
\node[box,draw=omExo,fill=omExo!8] (ship) at (6.4,0) {méthanier\\traversée};
\node[box,draw=omProp,fill=omProp!8] (reg) at (9.6,0) {regazéification\\terminal};
\draw[arr] (hh) -- (liq);
\draw[arr] (liq) -- (ship);
\draw[arr] (ship) -- (reg);
\node[font=\footnotesize,align=center,anchor=north] at (3.2,-0.75) {$1.15\times$ Henry Hub\\+ droit fixe};
\node[font=\footnotesize,align=center,anchor=north] at (6.4,-0.75) {affrètement, carburant,\\boil-off};
\node[font=\footnotesize,align=center,anchor=north] at (9.6,-0.75) {frais de regazéification~;\\vendu au hub};
\end{tikzpicture}

\omcaption{La chaîne GNL d'une cargaison américaine d'exportation, avec le coût que chaque maillon ajoute. Le \omterm{def:m3:physical-commodity-markets:netback}{netback}
est le prix de destination diminué des coûts situés à droite~; la cargaison est enlevée si le \omterm{def:m3:physical-commodity-markets:netback}{netback}
couvre le prix variable à la liquéfaction. Schéma de principe.}\label{fig:m3:natural-gas-and-lng:chain}
\end{omfigure}

\begin{example}[Europe ou Asie]\label{ex:m3:natural-gas-and-lng:choice}
Prenons une cargaison de 3.5 millions de MMBtu, un taux d'affrètement de 60\,000~\$ par jour, un boil-off de
0.1~\% par jour, et des traversées simples de 14 jours vers le nord-ouest de l'Europe et de 25 jours vers le nord de l'Asie
(à titre illustratif, mais proche des 14 jours de Corpus Christi à Rotterdam et des 28 jours jusqu'à Hong Kong
par le canal de Panama selon l'Oxford Institute for Energy Studies), avec des coûts de regazéification de \qty{0.50}{} et \qty{0.40}{\$/\mmbtu}. Avec
\omterm{def:m3:natural-gas-and-lng:hubs}{Henry Hub} à \qty{3.00}{}, le prix variable est de \qty{3.45}{\$/\mmbtu}~; le transport coûte
\qty{0.54}{} vers l'Europe et \qty{0.97}{\$/\mmbtu} vers l'Asie, de sorte que l'Asie doit payer
\qty{0.33}{\$/\mmbtu} de plus que l'Europe pour remporter la cargaison. À \qty{12.00}{} en Europe et
\qty{12.50}{} en Asie, les \omterm{def:m3:physical-commodity-markets:netback}{netbacks} sont de \qty{10.96}{} et \qty{11.13}{\$/\mmbtu}~: l'Asie.
\end{example}

\begin{omfigure}
\begin{tikzpicture}
\begin{axis}[width=11cm,height=5cm,ymode=log,
  ylabel={\$/MMBtu (échelle logarithmique)},
  tick label style={font=\small},label style={font=\small},
  xmin=2000,xmax=2026.75,ymin=1,ymax=100,grid=major,grid style={gray!25},
  ytick={2,5,10,20,50,100},yticklabels={2,5,10,20,50,100},
  xtick={2000,2004,2008,2012,2016,2020,2024},xticklabel style={/pgf/number format/1000 sep={}},
  legend columns=3,legend style={font=\footnotesize,at={(0.5,-0.18)},anchor=north,draw=none,fill=none,/tikz/every even column/.append style={column sep=8pt}},
  table/col sep=comma]
\addplot[omDef,thick] table[x=t,y=hh]{figdata/markets-3/04-natural-gas-and-lng/hubs.csv};
\addplot[omThm,thick] table[x=t,y=eu]{figdata/markets-3/04-natural-gas-and-lng/hubs.csv};
\addplot[omProp,thick,dashed] table[x=t,y=jp]{figdata/markets-3/04-natural-gas-and-lng/hubs.csv};
\legend{Henry Hub,hub européen,GNL asiatique}
\end{axis}
\end{tikzpicture}

\omcaption{Moyennes mensuelles de trois indices de référence gaziers, de janvier 2000 à juillet 2026, sur une
échelle logarithmique. Le gaz américain s'est découplé du reste du monde après 2009~; l'Europe a atteint un sommet de
\qty{69.98}{\$/\mmbtu} en août 2022. Données~: série FRED MHHNGSP (US Energy Information
Administration), PNGASEUUSDM et PNGASJPUSDM (International Monetary
Fund).}\label{fig:m3:natural-gas-and-lng:hubs}
\end{omfigure}

Appliquée mois par mois depuis les premières cargaisons d'exportation du Golfe américain en 2016, la règle de \omterm{def:m3:physical-commodity-markets:netback}{netback} de
l'\cref{ex:m3:natural-gas-and-lng:choice} envoie la cargaison vers l'Asie dans 93 mois sur 125 et
indique qu'elle ne devrait pas être enlevée sept mois en 2020 (de février à août) et en
février 2021, lorsqu'un gel au Texas a fait grimper \omterm{def:m3:natural-gas-and-lng:hubs}{Henry Hub}
(\cref{fig:m3:natural-gas-and-lng:netback}). Des acheteurs ont bien annulé des cargaisons en 2020~: la
US Energy Information Administration en a compté 45 annulées pour août et environ 30 pour
septembre, les prix spot bas en Europe et en Asie ayant rendu les exportations américaines non rentables.

\begin{omfigure}
\begin{tikzpicture}
\begin{axis}[width=11cm,height=4.8cm,ymode=log,
  ylabel={\$/MMBtu (échelle logarithmique)},
  tick label style={font=\small},label style={font=\small},
  xmin=2016,xmax=2026.75,ymin=0.5,ymax=100,grid=major,grid style={gray!25},
  ytick={1,2,5,10,20,50,100},yticklabels={1,2,5,10,20,50,100},
  xtick={2016,2018,2020,2022,2024,2026},xticklabel style={/pgf/number format/1000 sep={}},
  legend columns=2,legend style={font=\footnotesize,at={(0.5,-0.18)},anchor=north,draw=none,fill=none,/tikz/every even column/.append style={column sep=8pt}},
  table/col sep=comma]
\addplot[omDef,very thick] table[x=t,y=netback]{figdata/markets-3/04-natural-gas-and-lng/netback.csv};
\addplot[omThm,thick,dashed] table[x=t,y=fobvar]{figdata/markets-3/04-natural-gas-and-lng/netback.csv};
\legend{meilleur netback au Golfe américain,prix variable $1.15\times$ Henry Hub}
\end{axis}
\end{tikzpicture}

\omcaption{Meilleur \omterm{def:m3:physical-commodity-markets:netback}{netback} d'une cargaison du Golfe américain face à son prix variable, mensuel, de mars 2016
à juillet 2026. Lorsque le \omterm{def:m3:physical-commodity-markets:netback}{netback} passe sous la ligne en pointillés, la cargaison vaut
davantage annulée. Coûts d'expédition et de regazéification indicatifs de
l'\cref{ex:m3:natural-gas-and-lng:choice}~; prix tels que dans la
\cref{fig:m3:natural-gas-and-lng:hubs}.}\label{fig:m3:natural-gas-and-lng:netback}
\end{omfigure}

\begin{dated}{2026-09}{Exportations américaines}\label{dat:m3:natural-gas-and-lng:us}
Les exportations américaines de GNL ont augmenté de 26~\% pour atteindre 15.1 milliards de pieds cubes par jour en 2025, soit 26~\% du
total mondial~; les États-Unis sont le premier exportateur de GNL, devant le Qatar et l'Australie.
\end{dated}

\section{La crise de 2022}

Les exportations de gaz russe par gazoduc vers l'Union européenne et le Royaume-Uni ont chuté de presque
40~\% au cours des sept premiers mois de 2022 par rapport à 2021~; Nord Stream~1 a fonctionné à 20~\% de sa
capacité en juillet, puis s'est arrêté complètement au début de septembre (en novembre, le service de sécurité suédois
a établi que les explosions survenues sur les deux gazoducs Nord Stream relevaient d'un sabotage caractérisé)~; en octobre,
les livraisons russes par gazoduc étaient inférieures de 80~\% à leur niveau d'un an plus tôt. Les cours des hubs européens ont dépassé 300~euros le mégawattheure, l'échéance la plus proche culminant en août.
L'Europe a remplacé le gaz de gazoduc perdu par du GNL, en surenchérissant sur l'Asie pour les cargaisons flexibles,
et a rempli ses stocks selon les nouvelles règles. Les fournisseurs d'énergie qui avaient vendu leur production à terme sur les bourses ont dû faire face à des appels de marge se chiffrant en milliards~:
l'importateur Uniper avait tiré en totalité une ligne de crédit de 9~milliards d'euros accordée par l'État dès la fin août
(voir le \cref{ch:m3:when-markets-break-together}).

\begin{dated}{2026-09}{Le mécanisme de correction du marché}\label{dat:m3:natural-gas-and-lng:mcm}
La Council Regulation (EU) 2022/2578 a fixé une limite dynamique d'enchères sur les dérivés TTF
à échéance la plus proche~: lorsque le règlement-livraison de celle-ci dépassait 180~euros par MWh et se situait 35~euros au-dessus
d'un prix de référence mondial du GNL, les ordres supérieurs de plus de 35~euros par MWh à ce prix de référence ne devaient pas être
acceptés. Il s'appliquait à partir du 1\up{er} février 2023 et, prorogé une fois, jusqu'au 31 janvier 2025.
\end{dated}

\section{Tutoriel~: netbacks et choix de destination}\label{tut:m3:natural-gas-and-lng:netback}

\begin{tutorial}
\textbf{Objectif.} Calculer les \omterm{def:m3:physical-commodity-markets:netback}{netbacks} d'une cargaison, choisir sa destination et décider de la prendre en charge,
mois après mois. \textbf{État final~:}
\cref{ex:m3:natural-gas-and-lng:choice}, \cref{fig:m3:natural-gas-and-lng:netback} et
le nombre de mois sans prise en charge.
\begin{tutsteps}
\item \textbf{Expédition, \omterm{def:m3:physical-commodity-markets:netback}{netback}, choix.}
\omcode{code/firm/lngarb/firm_lngarb.py}{38}{57}{Coût d'expédition par MMBtu livré, netback et meilleure destination.}\label{lst:m3:natural-gas-and-lng:netback}
\item \textbf{Seuil de rentabilité et prise en charge.}
\omcode{code/firm/lngarb/firm_lngarb.py}{60}{68}{Spread de seuil de rentabilité entre deux destinations et règle de prise en charge.}\label{lst:m3:natural-gas-and-lng:lift}
\item \textbf{Exécuter} \texttt{m3\_lng.history()}, \texttt{summary()} et \texttt{fig\_lng.py}.
\end{tutsteps}
\textbf{Pour aller plus loin.} Faire passer les cargaisons asiatiques par le cap de Bonne-Espérance lorsque le canal de
Panama est restreint, et compter combien de décisions de destination changent~; faire varier le taux d'affrètement avec le
spread, comme il le fait lorsque de nombreuses cargaisons se disputent le même arbitrage.
\end{tutorial}

\section{Construction~: le moteur d'arbitrage des cargaisons}\label{bld:m3:natural-gas-and-lng:lngarb}

\begin{build}
\bfield{Rôle} La société modèle dispose d'une capacité de liquéfaction~: chaque mois, elle doit
décider s'il faut prendre en charge chaque cargaison et où l'envoyer, et elle cote le gaz sur les hubs dans
trois unités.
\bfield{Interface} \texttt{eur\_mwh\_to\_usd\_mmbtu}, \texttt{usd\_mmbtu\_to\_eur\_mwh}~;
\texttt{fob\_price(henry\_hub, slope, fee)}~; \texttt{Route(name, days, charter\_per\_day,
boiloff\_per\_day, regas)}~; \texttt{shipping\_cost}~; \texttt{netback}~; \texttt{choose}~;
\texttt{breakeven\_spread}~; \texttt{lift}.
\bfield{Règles} Les frais fixes sont irrécupérables et n'entrent jamais dans la décision de prise en charge~; l'évaporation
est valorisée à la valeur FOB de la cargaison et réduit le volume livré~; l'affrètement est payé pour l'aller-retour~; en cas
d'égalité, la destination venant en premier dans l'ordre alphabétique l'emporte.
\bfield{Tests d'acceptation} \texttt{code/firm/lngarb/tests/}~: allers-retours unitaires~; une route plus longue
coûte davantage~; le choix bascule au spread de seuil de rentabilité~; la règle de prise en charge.
\bfield{Pour aller plus loin} Contraintes de canal et temps d'attente qui en résultent~; une flotte de transporteurs
programmée sur une année~; la valeur de l'option de prise en charge calculée par Monte-Carlo sur des prix de hubs
corrélés.
\end{build}

\begin{omsources}
\item Regulation (EU) 2022/1032 et Regulation (EU) 2025/1733 sur le stockage de gaz~; Council
Regulation (EU) 2022/2578 (mécanisme de correction du marché)~; Commission Regulation (EU)
2015/703 (jour gazier).
\item US Energy Information Administration, \emph{Today in Energy}~: exportations de GNL en 2020 et
2025~; exportations russes par gazoduc vers l'Europe (2022). IEA, \emph{2022 Energy Crisis: Frequently
Asked Questions}.
\item Cheniere Energy Partners, Form 10-K (tarification SPA à 115~\% de Henry Hub plus des
frais fixes).
\item Swedish Security Service, ``Confirmed sabotage of the Nord Stream gas pipelines'', 18 novembre 2022 (copie
sur Internet Archive)~; Fortum, communiqué de presse sur la facilité KfW d'Uniper, 29 août 2022~; Oxford Institute for Energy Studies, \emph{LNG
Shipping Chokepoints}, NG 188, février 2024~; ACER, rapport préliminaire sur le
mécanisme de correction du marché, 2023.
\item CME Group, Henry Hub natural gas futures overview~; S\&P Global Platts, JKM.
\item D.~Dumitrescu et al.,
``European energy crisis, the Dutch TTF, and the market correction mechanism'', \emph{Journal of World Energy Law \& Business}, 2024.
\item Séries FRED MHHNGSP, PNGASEUUSDM, PNGASJPUSDM.
\end{omsources}

\section{Exercices}

\begin{exercise}[$\star$]\label{exo:m3:natural-gas-and-lng:1}
Convertir \qty{10}{\$/\mmbtu} en EUR/MWh au taux EURUSD de 1.16, et en pence par therm au
taux GBPUSD de 1.30.
\end{exercise}

\begin{exercise}[$\star$]\label{exo:m3:natural-gas-and-lng:2}
Pourquoi le \omterm{def:m3:natural-gas-and-lng:season}{spread été-hiver} est-il positif la plupart des années, et qu'est-ce qui le borne~?
\end{exercise}

\begin{exercise}[$\star$]\label{exo:m3:natural-gas-and-lng:3}
Un acheteur américain de GNL paie des frais fixes de \qty{2.50}{\$/\mmbtu} et 115~\% de \omterm{def:m3:natural-gas-and-lng:hubs}{Henry Hub}. \omterm{def:m3:natural-gas-and-lng:hubs}{Henry Hub}
vaut \qty{2.00}{} et son meilleur \omterm{def:m3:physical-commodity-markets:netback}{netback} \qty{2.10}{\$/\mmbtu}. Prend-il la cargaison en charge~? Que perd-il
par MMBtu de capacité dans les deux cas~?
\end{exercise}

\begin{exercise}[$\star\star$]\label{exo:m3:natural-gas-and-lng:4}
Avec les routes de l'\cref{ex:m3:natural-gas-and-lng:choice} et \omterm{def:m3:natural-gas-and-lng:hubs}{Henry Hub} à \$3, l'Europe
paie \qty{11.00}{\$/\mmbtu}. Au-dessus de quel prix asiatique la cargaison part-elle pour l'Asie~?
\end{exercise}

\begin{exercise}[$\star\star$]\label{exo:m3:natural-gas-and-lng:5}
Un hub dont les gazoducs sortants sont saturés se traite \qty{2}{\$/\mmbtu} sous \omterm{def:m3:natural-gas-and-lng:hubs}{Henry Hub} alors que
le coût de transport vaut \qty{0.30}{}. Qui empoche l'écart, et quelle transaction conviendrait-il de réaliser si
l'on détient des capacités~?
\end{exercise}

\begin{exercise}[$\star\star$]\label{exo:m3:natural-gas-and-lng:6}
Le voyage vers l'Asie double et passe à 50 jours (un canal est fermé). De combien le spread du point mort
augmente-t-il, les autres données étant celles de l'\cref{ex:m3:natural-gas-and-lng:choice}~?
\end{exercise}

\begin{exercise}[$\star\star\star$]\label{exo:m3:natural-gas-and-lng:7}
\emph{Programmation.} Avec \texttt{history}, comptez les mois depuis mars 2016 où
la cargaison part vers l'Asie et listez ceux où elle n'est pas enlevée.
\end{exercise}

\begin{exercise}[$\star\star\star$]\label{exo:m3:natural-gas-and-lng:8}
\emph{Repérer l'erreur.} «~Le gaz européen était dix fois plus cher que le gaz américain en août 2022, donc
chaque cargaison américaine rapportait dix fois son coût.~»
\end{exercise}

\section{Problème~: Où envoyer la cargaison~?}

\begin{problem}[{Problème du week-end --- une cargaison flexible entre deux continents}]\label{pb:m3:natural-gas-and-lng:1}
Une \omterm{def:m3:physical-commodity-markets:house}{maison de négoce} détient une cargaison du golfe du Mexique de 3.5~millions de MMBtu, avec une
\omterm{def:m3:natural-gas-and-lng:contracts}{flexibilité de destination} complète, achetée à 115~\% de \omterm{def:m3:natural-gas-and-lng:hubs}{Henry Hub} plus des frais irrécupérables. Utilisez les routes et les coûts de
l'\cref{ex:m3:natural-gas-and-lng:choice}.

\medskip\noindent\textbf{Partie I --- Les coûts.}
\begin{enumerate}
\item \omterm{def:m3:natural-gas-and-lng:hubs}{Henry Hub} valant 3~\$, donnez le prix variable de la cargaison.
\item Donnez le coût de transport par MMBtu livré en Europe.
\item Et en Asie.
\item Quelle part de la cargaison s'évapore à chaque voyage~?
\item Pourquoi l'affrètement est-il payé sur deux fois la durée du trajet simple~?
\end{enumerate}

\medskip\noindent\textbf{Partie II --- La décision.}
\begin{enumerate}[resume]
\item Donnez le spread du point mort, défini comme l'Asie moins l'Europe.
\item L'Europe à 12.00~\$ et l'Asie à 12.50~\$, où part la cargaison~?
\item Quelle est la marge de la cargaison au-dessus de son prix variable, en dollars~?
\item Que devient la décision si le taux d'affrètement double~?
\item En août 2022 (\omterm{def:m3:natural-gas-and-lng:hubs}{Henry Hub} à 8.81~\$, Europe à 69.98~\$, Asie à 54.16~\$), où part-elle~?
\end{enumerate}

\medskip\noindent\textbf{Partie III --- L'option.}
\begin{enumerate}[resume]
\item Pourquoi le droit de choisir une destination a-t-il une valeur monétaire, même lorsque les prix sont égaux
en moyenne~?
\item En juin 2020, que disait la règle du \omterm{def:m3:physical-commodity-markets:netback}{netback}, et pourquoi~?
\item Que changent les frais irrécupérables dans cette décision~?
\item Qui y gagne lorsque l'acheteur annule une cargaison~?
\item Quel marché bouge lorsque de nombreuses cargaisons flexibles poursuivent le même arbitrage~?
\end{enumerate}

\medskip\noindent\textbf{Partie IV --- Jugement.}
\begin{enumerate}[resume]
\item Pourquoi l'Europe a-t-elle remporté la plupart des cargaisons flexibles en 2022~?
\item Pourquoi la fermeture d'un canal fait-elle monter le prix du gaz en Asie par rapport à l'Europe~?
\item Pourquoi un acheteur signerait-il un contrat indexé sur le pétrole plutôt qu'un contrat lié à \omterm{def:m3:natural-gas-and-lng:hubs}{Henry Hub}~?
\item Énoncez le \emph{résultat nommé}~: le spread auquel la cargaison est indifférente entre
les deux destinations.
\item En une phrase~: qu'est-ce qui relie les prix du gaz dans le monde~?
\end{enumerate}
\end{problem}

\section{Questions d'entretien}

\begin{interviewq}[$\star$ \iqroles{trader}]\label{iq:m3:natural-gas-and-lng:1}
Quelle est la différence entre \omterm{def:m3:natural-gas-and-lng:hubs}{Henry Hub} et le TTF en tant que hubs~?
\end{interviewq}

\begin{interviewq}[$\star$ \iqroles{trader, chercheur}]\label{iq:m3:natural-gas-and-lng:2}
Pourquoi le gaz est-il saisonnier alors que le pétrole l'est beaucoup moins~?
\end{interviewq}

\begin{interviewq}[$\star\star$ \iqroles{chercheur}]\label{iq:m3:natural-gas-and-lng:3}
Un contrat d'enlèvement de GNL américain est un strip d'options. Sur quoi, et avec quel strike~?
\end{interviewq}

\begin{interviewq}[$\star\star$ \iqroles{trader, risque}]\label{iq:m3:natural-gas-and-lng:4}
Vous êtes long de gaz sur un hub et short sur \omterm{def:m3:natural-gas-and-lng:hubs}{Henry Hub}. Qu'est-ce qui peut vous faire perdre beaucoup en une journée~?
\end{interviewq}

\begin{interviewq}[$\star\star$ \iqroles{risque}]\label{iq:m3:natural-gas-and-lng:5}
Pourquoi les fournisseurs d'énergie européens, pourtant entièrement couverts, ont-ils fait face à une crise de liquidité en 2022~?
\end{interviewq}

\begin{interviewq}[$\star\star\star$ \iqroles{développeur, trader}]\label{iq:m3:natural-gas-and-lng:6}
Concevez l'outil qui indique chaque matin à un desk GNL où doit aller chacune de ses vingt
cargaisons.
\end{interviewq}
